\section{讨论：可证伪边界、接口假设与可检验预测}

\subsection{分层可证伪结构}

本文将可证伪性分为三层：
\begin{itemize}
  \item \textbf{M 层（数学层）}：计数定理、折叠映射性质、循环/边界分裂、退化度分布等，完全由有限构造确定。
  \item \textbf{P 层（协议层）}：将扫描--投影读出解释为“观测协议”，并将稳定子空间条件化解释为“只能观测存在态”的形式化表达。该层的对象是协议等价类与可重复统计。
  \item \textbf{I 层（解释层）}：将 $\varph$--$\pi$--$e$ 通道对应到具体物理/信息学语义（如“闭合相位”“解析稳定域”“对数阻抗”）。该层不作为定理链条的前提，而作为可检验建模选择。
\end{itemize}

\subsection{可检验预测（示例）}

\begin{enumerate}
  \item \textbf{稳定类型数量的刚性增长律：}在局部禁止律为 $11$ 的最小语法下，窗口长度 $m$ 的稳定类型数严格为 $F_{m+2}$。
  \item \textbf{循环/边界分裂：}环接可容许性给出 $\card{X_m^{\mathrm{cyc}}}=F_{m-1}+F_{m+1}$、$\card{X_m^{\mathrm{bdry}}}=F_{m-2}$。
  \item \textbf{折叠压缩的可复核退化度：}$\mathrm{Fold}_m$ 的平均退化度为 $2^m/F_{m+2}$，并随 $m$ 以 $(2/\varph)^m$ 的指数尺度增长。
  \item \textbf{最小例 $m=6$ 的刚性结构：}$64\to 21$ 压缩、$3\oplus 18$ 分裂及边界双前像结构构成最小“黄金存在”实例，可通过穷举完全复核。
  \item \textbf{通道分离条件：}在更长窗口或加权语法中，解析稳定性（$e$ 通道）与语法合法性（$\varph$ 通道）可分离；该分离的出现与否及其尺度为模型扩展的直接证伪点。
\end{enumerate}

